\documentclass[10pt,a4paper]{article}
\usepackage[T1]{fontenc}
\usepackage[utf8]{inputenc}
\usepackage{lmodern}
\usepackage[margin=20mm]{geometry}
\usepackage{amsmath,amssymb,array,longtable,booktabs,calc,enumitem}
\usepackage{xurl,xcolor,mdframed}
\usepackage[unicode,hidelinks]{hyperref}
\setlength{\parindent}{0pt}\setlength{\parskip}{5pt}
\setlength{\emergencystretch}{3em}\setlength{\tabcolsep}{4pt}
\renewcommand{\arraystretch}{1.12}\allowdisplaybreaks
\newcolumntype{L}[1]{>{\raggedright\arraybackslash}p{#1}}
\newcommand{\proj}{\operatorname{proj}}
\newcommand{\fixedcosts}{\sum_c F_cy_c+\sum_c K_cq_c}
\newcommand{\note}{\textbf{Note.}\ }
\begin{document}
\section{ER-DRO model (with CVaR-type ambiguity set)}
Data: $(y^s,x^s)$, $s\in S$---for example, city-year observations of choice probabilities $y^s$ and city characteristics $x^s$. Each observation corresponds to one scenario, so $S$ is the set of observation indices. Regression gives $\hat f$ and residuals $\hat\varepsilon^s=y^s-\hat f(x^s)$ (4.A); the characteristics $x_c$ of candidate city $c$ are used to generate scenarios $D_c^s$ (4.B).
{\small
\[
\boxed{\begin{aligned}
\min_{y,q,O,U,\tau,w}\quad
&\underbrace{\sum_cF_cy_c}_{\text{Entry costs}}+
\underbrace{\sum_cK_cq_c}_{\text{Vehicle costs}}+
\underbrace{\tau+\frac{1}{|S|(1-\zeta)}\sum_{s\in S}w^s}_{\text{CVaR of scenario penalties}}
\end{aligned}}
\]}
\note The baseline term $\frac1{|S|}\sum_s H(q,D^s)$ is replaced by $\tau+\frac1{|S|(1-\zeta)}\sum_sw^s$, where $H(q,D^s)=c_c^+O_c^s+c_c^-U_c^s$. Context $x_c$ enters only the scenario values $D_c^s=M_cp_c^s$; the optimization structure remains unchanged.

\section{Constraints}
\[
\boxed{\begin{aligned}
w^s&\ge\sum_c\left(c_c^+O_c^s+c_c^-U_c^s\right)-\tau &&\forall s\in S,\\
\kappa_cq_c-D_c^s&=O_c^s-U_c^s &&\forall c\in C,\ s\in S,\\
0\le q_c&\le\bar q_cy_c &&\forall c\in C,\\
\sum_{c\in C}q_c&\le Q,\\
y_c\in\{0,1\},\quad q_c&\in\mathbb Z_{\ge0} &&\forall c\in C,\\
O_c^s,U_c^s&\ge0 &&\forall c\in C,\ s\in S,\\
w^s&\ge0 &&\forall s\in S,\qquad\tau\in\mathbb R.
\end{aligned}}
\]
{\footnotesize
\begin{longtable}{@{}L{0.40\textwidth}L{0.56\textwidth}@{}}
\toprule Constraint & Meaning\\\midrule\endhead\bottomrule\endlastfoot
$0\le q_c\le\bar q_cy_c$ & Zero vehicles if the city is not entered; up to the city-specific limit if it is entered\\
$\sum_cq_c\le Q$ & Total vehicle limit\\
$\kappa_cq_c-D_c^s=O_c^s-U_c^s$ & The difference between service capacity and robotaxi demand equals unused capacity minus unmet demand.\\
\textcolor{blue}{$w^s\ge\sum_c(c_c^+O_c^s+c_c^-U_c^s)-\tau$, $w^s\ge0$} & Excess of scenario $s$'s penalty over threshold $\tau$. Optimization makes $w^s$ as small as possible, so \textcolor{blue}{$w^s=(\text{penalty}_s-\tau)^+$}.\\
\end{longtable}}

\section{Notation}
\subsection{\textbf{Contextual + Additional DRO Notation}}
{\footnotesize
\begin{longtable}{@{}L{0.14\textwidth}L{0.45\textwidth}L{0.35\textwidth}@{}}
\toprule Symbol & Meaning & Type and unit\\\midrule\endhead\bottomrule\endlastfoot
$(y^s,x^s)$ & Dependent variable (choice probability) and covariates of observation $s$ & Data; $s\in S$\\
$x_c$ & Characteristic (covariate) vector of candidate city $c$ & Input; corresponds to Kannan's new observation $x$\\
$\hat f$ & Estimated regression function; $\hat f(x)\approx\mathbb E[y\mid x]$ & Estimated function\\
$\hat\varepsilon^s$ & Residual $y^s-\hat f(x^s)$ & Derived; on the $p$ scale\\
$\zeta$ & Radius (tail proportion) of the CVaR-type ambiguity set. $1-\zeta$ is the share of scenarios retained as the worst group & Input; $[0,1)$\\
$\tau$ & Penalty threshold. A decision variable; at the optimum, the ``$|S|(1-\zeta)$-th largest penalty'' & Decision variable; currency\\
$w^s$ & Amount by which scenario $s$'s penalty exceeds $\tau$ & Decision variable; currency\\
\end{longtable}}
\subsection{Baseline Notation}
{\footnotesize
\begin{longtable}{@{}L{0.10\textwidth}L{0.43\textwidth}L{0.41\textwidth}@{}}
\toprule Symbol & Meaning & Type and unit\\\midrule\endfirsthead
\toprule Symbol & Meaning & Type and unit\\\midrule\endhead\bottomrule\endlastfoot
$C$ & Set of potential candidate cities & Input set; cities\\
$S$ & Set of demand scenarios & Input set; scenarios. Here it is the set of regression-data observation indices; $|S|=$ number of observations\\
$y_c$ & 1 if city $c$ is entered, else 0 & Decision variable; binary (0 or 1)\\
$q_c$ & Number of vehicles deployed in city $c$ & Decision variable; vehicles (integer)\\
$\kappa_c$ & Number of trips one vehicle can serve during the period & Scale parameter; trips/vehicle over the evaluation period\\
$Q$ & Total number of vehicles available for deployment & Input parameter; vehicles\\
$\bar q_c$ & Upper bound on vehicle deployment in each city & Input parameter; vehicles\\
$F_c$ & Fixed cost of city entry and basic infrastructure & Input parameter; currency over the evaluation period\\
$K_c$ & Per-vehicle preparation, deployment, and holding cost over the evaluation period & Input parameter; currency/vehicle over the evaluation period\\
$c_c^+$ & Additional loss per unit of unused service capacity & Input parameter; currency/trip of unused capacity\\
$c_c^-$ & Foregone contribution margin per unit of unserved potential demand & Input parameter; currency/unserved trip\\
$M_c$ & Total number of \textbf{potential travel opportunities across all modes} in city $c$ during one year & Input parameter; usage opportunities over the evaluation period. Data or assumption\\
\textcolor{blue}{$p_c^s$} & \textcolor{blue}{Probability of choosing robotaxi service for a travel opportunity in city $c$} & \textcolor{blue}{Derived input (generated by regression); dimensionless probability in $[0,1]$. $p_c^s=\proj_{[0,1]}(\hat f(x_c)+\hat\varepsilon^s)$}\\
$D_c^s$ & Robotaxi choice demand & Derived input; trips over the evaluation period. $D_c^s=M_cp_c^s$\\
$O_c^s$ & Unused service capacity & Outcome variable; trips over the evaluation period. $O_c^s=(\kappa_cq_c-D_c^s)^+$\\
$U_c^s$ & Unserved robotaxi demand & Outcome variable; trips over the evaluation period. $U_c^s=(D_c^s-\kappa_cq_c)^+$\\
\end{longtable}}

\section{Explanations}
\subsection{Derivation of the CVaR Formulation with Contextual Demand}
\textbf{A. Definitions.}
\begin{itemize}
\item $H(q,D^s)$: scenario $s$'s second-stage cost (``Common second stage'' at the top of the document).
\[
\begin{aligned}
H(q,D^s)&=\min_{O^s,U^s\ge0}\sum_c(c_c^+O_c^s+c_c^-U_c^s)\\
&\text{s.t.}\quad\kappa_cq_c-D_c^s=O_c^s-U_c^s\quad\forall c.
\end{aligned}
\]
At the optimum, $O_c^s=(\kappa_cq_c-D_c^s)^+$ and $U_c^s=(D_c^s-\kappa_cq_c)^+$, so
\[
H(q,D^s)=\sum_c\left[c_c^+(\kappa_cq_c-D_c^s)^++c_c^-(D_c^s-\kappa_cq_c)^+\right].
\]
\item $w^s$: the amount by which scenario $s$'s penalty exceeds threshold $\tau$. At the optimum, $w^s=(H(q,D^s)-\tau)^+$.
\item $D_c^s(x_c)=D_c^s=M_cp_c^s=M_c\proj_{[0,1]}(\hat f(x_c)+\hat\varepsilon^s)$,\quad $D^s=(D_1^s,D_2^s,\ldots,D_{|C|}^s)$.
\item
\[
\hat P_n^{\rm ER}(x)=\frac1{|S|}\sum_{s\in S}\delta_{D^s(x)},\qquad
\hat{\mathcal P}_n(x)=\left\{\sum_sp_s\delta_{D^s(x)}:p\in\mathcal P_n(x;\zeta)\right\}.
\]
\end{itemize}
\textbf{B. Chain of equalities.} Starting from Kannan's residuals-based DRO formulation
\[
\min_{z\in Z}\left[\sup_{Q\in\hat{\mathcal P}_n(x)}\mathbb E_{Y\sim Q}[c(z,Y)]\right],
\]
fix the decision $z=(y,q)$ and expand the inner supremum first.
{\small
\begin{align*}
&\sup_{Q\in\hat{\mathcal P}_n(x)}\mathbb E_{Y\sim Q}[c(z,Y)]\\
&\overset{(a)}{=}\sup_{p\in\mathcal P_n(x;\zeta)}\sum_{s\in S}p_sc(z,D^s(x))\\
&\overset{(b)}{=}\fixedcosts+\sup_{p\in\mathcal P_n}\sum_{s\in S}p_sH(q,D^s)\\
&\overset{(c)}{=}\fixedcosts+\min_{\tau\in\mathbb R}\left[\tau+\frac1{|S|(1-\zeta)}\sum_{s\in S}(H(q,D^s)-\tau)^+\right]\\
&\overset{(d)}{=}\fixedcosts+\min_{\tau,w}\left\{\begin{aligned}
&\tau+\frac1{|S|(1-\zeta)}\sum_{s\in S}w^s:\\
&w^s\ge H(q,D^s)-\tau,\quad w^s\ge0\quad\forall s
\end{aligned}\right\}\\
&\overset{(e)}{=}\fixedcosts+\min_{\tau,w,O,U}\left\{\begin{aligned}
&\tau+\frac1{|S|(1-\zeta)}\sum_{s\in S}w^s:\\
&w^s\ge\sum_c(c_c^+O_c^s+c_c^-U_c^s)-\tau,\\
&\kappa_cq_c-D_c^s=O_c^s-U_c^s,\\
&w^s,O_c^s,U_c^s\ge0\quad\forall c,s
\end{aligned}\right\}.
\end{align*}}
Finally, apply the outer $\min_{y,q}$. Since $\min_{y,q}\min_{\tau,w,O,U}=\min_{y,q,O,U,\tau,w}$, the result is
\[
\begin{aligned}
\text{(f)}\quad\min_{y,q,O,U,\tau,w}\quad&\fixedcosts+\tau+\frac1{|S|(1-\zeta)}\sum_{s\in S}w^s\\
\text{s.t.}\quad&w^s\ge\sum_c(c_c^+O_c^s+c_c^-U_c^s)-\tau,\quad w^s\ge0\quad\forall s\in S,\\
&\kappa_cq_c-D_c^s=O_c^s-U_c^s,\quad O_c^s,U_c^s\ge0\quad\forall c\in C,\ s\in S,\\
&0\le q_c\le\bar q_cy_c,\quad\sum_cq_c\le Q,\\
&y_c\in\{0,1\},\quad q_c\in\mathbb Z_{\ge0},\quad\tau\in\mathbb R.
\end{aligned}
\]
\textbf{C. Why each equality holds.}
\begin{itemize}
\item (a) The ambiguity set has the same support, so $Q=\sum_sp_s\delta_{D^s}$ and expectation is a weighted sum.
\item (b) In $c(z,D^s)=\sum_cF_cy_c+\sum_cK_cq_c+H(q,D^s)$, the first two terms do not depend on $s$. Since $\sum_sp_s=1$, they can be taken outside the supremum.
\item (c) If $\mathcal P_n$ is the box in Example 1, $p_s\le1/(|S|(1-\zeta))$, the supremum is the average of the $|S|(1-\zeta)$ largest penalties, namely $\operatorname{CVaR}_\zeta$. This is equivalent to the Rockafellar--Uryasev threshold representation.
\item (d) Linearization of $(\cdot)^+$. Minimizing $w^s$ under $w^s\ge H-\tau$ and $w^s\ge0$ gives $w^s=(H-\tau)^+$, so the values agree.
\item (e) $H(q,D^s)$ is itself a minimum over $O^s,U^s$. ``$w^s\ge\min_{O,U}(\cdots)-\tau$'' is equivalent to ``$w^s\ge(\cdots)-\tau$ for some feasible $O^s,U^s$.'' The inner minimum can therefore be merged into the outer minimum and the balance equation imposed as a constraint.
\item (f) The outer $\min_{y,q}$ and inner $\min_{\tau,w,O,U}$ are both minimizations and can be combined. The supremum has already become a minimum in (c), so no min--max structure remains.
\end{itemize}

\subsection{Scenario Generation}
\textbf{A. Regression}

For the characteristics $x_c$ of candidate city $c$,
{\color{blue}\[
p_c^s=\proj_{[0,1]}(\hat f(x_c)+\hat\varepsilon^s),\qquad D_c^s=M_cp_c^s,\qquad s\in S.
\]}
Fit $\hat f$ and calculate residuals $\hat\varepsilon^s=y^s-\hat f(x^s)$, $s\in S$. $x^s$ are training points. \textcolor{red}{$x$ and $\epsilon$ are assumed to follow $N(0,1)$ and to be independent of the covariate $x$ (homoscedasticity must be checked after regression).} As in the baseline model, $M_c$ is held fixed.

\textcolor{blue}{\textbf{B. What are the covariates $X$?}}

\textbf{B-1. How is $M$ determined?} For details, see \texttt{2.driving factors/C.Selection/Factors\_of\_M}.

$M_c$ is constructed from \textcolor{blue}{NHTS group-specific per-person trip rates $\times$ city populations}. This gives a fixed value calculated for all candidate cities using the same definition, without regression.

\textbf{B-2. How is $p$ determined?}

For details, see \texttt{2.driving factors/C.Selection/Factors\_of\_p}.

The preliminary selection identifies two city characteristics and three decision-dependent service variables as initial priorities. Test the two exogenous variables in the current model first; endogenous service variables will be incorporated into a decision-dependent ER-DRO extension. Exogenous/endogenous refers to dependence on operating decisions.
\begin{itemize}
\item \textcolor{blue}{\textbf{Exogenous:} zero-vehicle household share (exogenous), age composition (exogenous).}
\item \textcolor{blue}{\textbf{Endogenous:} relative fare (endogenous), relative waiting time (endogenous), in-vehicle support features (endogenous).}
\end{itemize}
\textcolor{red}{\textbf{C. Which regression method will be used, and why? Still to be decided}}
\begin{mdframed}[linewidth=0.6pt]
\footnotesize
\note \textbf{Kannan et al. (Residuals-based Distributionally Robust Optimization)}

Can any regression method be used? The paper's framework is method-agnostic (OLS, Lasso, Ridge, kNN, and random forests are all possible). Three considerations nevertheless matter.
\begin{itemize}
\item \textbf{Theoretical conditions.} Assumptions 2, 8, and 15 concern how quickly regression error decreases with $|S|$. Parametric methods such as OLS and Lasso satisfy $r=1$ (error $\sim|S|^{-1/2}$); nonparametric methods are slower, with $r=O(1)/d_x$. With few cities and many features, parametric methods are needed for the theory to apply.
\item \textbf{Regularize with small samples.} In the paper's experiments, Lasso and Ridge outperformed OLS when $|S|$ was small relative to $d_x$ and complemented DRO. Small $|S|$ also makes training residuals optimistically small, so leave-one-out (jackknife) residuals from the previous paper [34] are preferable.
\item \textbf{Model-structure assumption.} In $Y=f^*(X)+\epsilon$, errors must be independent of $X$ (homoscedastic) for residuals to be shared across cities. Regressing trip counts $D$ violates this assumption because larger cities tend to have larger errors. Regressing choice probability $p$ (or $\operatorname{logit}p$) on city characteristics and constructing $D_c=M_cp_c$ is much closer to homoscedasticity; projection onto $[0,1]$ is also natural. If heteroscedasticity remains, use the authors' subsequent extension [35].
\end{itemize}
\end{mdframed}

\textbf{D. Selecting the Ambiguity Set}

We use the CVaR-based ambiguity sets with the same support as $\hat P_n^{\rm ER}$, the empirical distribution of the residuals.
\begin{mdframed}[linewidth=0.6pt]
\footnotesize
\note \textbf{Kannan et al. (Residuals-based Distributionally Robust Optimization)}

Ambiguity sets with the same support as $\hat P_n^{\rm ER}(x)$ (cf. [4, 6], for instance): given radius $\zeta_n(x)\ge0$, set
\[
\hat{\mathcal P}_n(x)=\left\{Q=\sum_{i=1}^np_i\delta_{\proj_Y(\hat f_n(x)+\hat\varepsilon_n^i)}:p\in\mathcal P_n(x;\zeta_n(x))\right\},
\]
where $\mathcal P_n(x;\zeta_n(x))$ is a generic ambiguity set for the $n$-dimensional vector of probabilities $p$. We focus on sets $\mathcal P_n(x;\zeta_n(x))$ that satisfy for each $x\in X$
\[
p\in\mathbb R_+^n\quad\text{and}\quad\sum_{i=1}^np_i=1,\quad\forall p\in\mathcal P_n(x;\zeta_n(x)),
\]
\begin{equation*}
\lim_{\zeta\downarrow0}\mathcal P_n(x;\zeta)=\mathcal P_n(x;0)=\left\{\left(\frac1n,\frac1n,\ldots,\frac1n\right)\right\}.\tag{8}
\end{equation*}
\textbf{Example 1.} CVaR-based ambiguity set [45, 48]: given radius $\zeta_n(x)\in[0,1)$, set
\[
\mathcal P_n(x;\zeta_n(x)):=\left\{p\in\mathbb R_+^n:\sum_{i=1}^np_i=1,\quad p_i\le\frac1{n(1-\zeta_n(x))},\quad\forall i\in[n]\right\}.
\]
Observe that $\zeta_n(x)$ enters the ambiguity set $\mathcal P_n(x;\zeta_n(x))$ through the CVaR risk parameter.
\end{mdframed}

\textbf{D-1. Why choose CVaR?}

If we do not fully trust assigning probability $1/|S|$ to each of the $|S|$ scenarios, two things may be questioned. Doubting scenario values leads to sets that move points (Wasserstein, SRO); doubting scenario probabilities leads to sets that change only weights (CVaR-type). We choose the latter.
\begin{itemize}
\item \textbf{The structure is preserved.} The ``average of the $(1-\zeta)$ fraction of scenarios with the largest penalties'' can be added directly to the existing MILP using only threshold $\tau$ and excess variables $w^s$. Point-moving sets introduce min--max--min within integer decisions; the paper also regards this as generally difficult.
\item \textbf{Point-moving sets have different meanings depending on the $c^+/c^-$ ratio.} If $c^+\ll c^-$, the worst-case distribution pushes demand only upward, making the experiment effectively similar to demand-multiplier sensitivity. As $c^+$ approaches $c^-$, it also pushes toward unused capacity, complicating interpretation. CVaR-type sets ask the same question regardless of the ratio: ``What if bad scenarios occur more frequently?'' This makes comparisons across ratio combinations clearer.
\item \textbf{Frequency has the weakest basis, rather than values.} ER-DRO scenario values are grounded in data---regression predictions + actual residuals---but with fewer observations, assigning exactly $1/|S|$ to every scenario is the most fragile assumption. A probability-only ambiguity set targets that weakness.
\item \textbf{Easy to explain.} $\zeta=0.5$ means ``minimize the average of the worst half.'' Results can be read through the ``value of one vehicle'' logic by replacing scenario proportions with worst-case weights.
\end{itemize}
Caveat: in the paper's experiments, the Hellinger-type set from the same class yielded smaller improvements than Wasserstein. Those small-sample continuous-distribution settings differ from ours, but validation is needed; run both sets and provide a comparison table.

\textbf{D-2. Understanding the ``$\tau$ trick''}

\note Example (the numbers below are from the v3 instance without covariates, i.e., ten manually designed scenarios).

Fix the reference allocation at $(132,368,500)$ and set $\zeta=0.5$ (worst five scenarios). Scenario penalties $C_s=\sum_cc_c^-U_c^s$ and the solver's $\tau,w^s$ are as follows (million USD).
{\footnotesize
\begin{center}
\begin{tabular}{@{}lrrrrrrrrrr@{}}
\toprule
$s$ & 8 & 4 & 5 & 6 & 10 & 9 & 1 & 3 & 2 & 7\\\midrule
$C_s$ & 297.7 & 288.0 & 220.6 & 174.4 & 156.5 & 114.7 & 104.0 & 58.8 & 56.9 & 6.3\\
$w^s=(C_s-\tau)^+$ & 141.2 & 131.5 & 64.1 & 17.9 & 0 & 0 & 0 & 0 & 0 & 0\\
\bottomrule
\end{tabular}
\end{center}}
$\tau=156.5$ (the fifth-largest penalty), and the objective is
\[
18.47+156.5+\frac15(141.2+131.5+64.1+17.9)=245.9.
\]
This is exactly fixed costs + the worst-five average, 227.4. Thus $w^s$ is ``the amount by which scenario $s$'s penalty exceeds the threshold,'' and is zero outside the worst group. Changing the allocation changes $C_s$, and $\tau,w^s$ change accordingly.

\textbf{I do not set $\tau$.} Like $y,q$, $\tau$ is a decision variable found by the solver. At the optimum it is always the ``$|S|(1-\zeta)$-th largest penalty.'' I set only $\zeta$ (using one of the three methods in 4.D).

\note \textbf{References}
\begin{itemize}
\item Rockafellar \& Uryasev (2000), ``Optimization of conditional value-at-risk,'' \emph{Journal of Risk}. Origin of the $\tau$ trick. They use $\alpha$ for the threshold and $\beta$ for the tail level (our $\tau,\zeta$).
\item Shapiro, Dentcheva, Ruszczy\'{n}ski, \emph{Lectures on Stochastic Programming}, chapter on risk-averse optimization. The equivalence of CVaR's ``$\min_\tau$ threshold representation'' and ``worst-case weights with probabilities capped at $1/(|S|(1-\zeta))$'' is presented there. The paper's Example 1 uses the latter directly.
\item Section 7 of Kannan's paper: the portfolio objective is written as $\min_\tau\mathbb E[\tau+\frac1{1-\beta}\max\{0,\cdot-\tau\}]$ and solved as an LP. This is the same approach as ours.
\end{itemize}

\textbf{E. Methods for Selecting $\zeta$ (Radius)}

In the CVaR-type set, $\zeta$ is not a distance radius but a tail proportion determining ``how many worst scenarios to consider'' ($1-\zeta=$ retained proportion).

\textbf{E-1. Cross validation}

\textbf{E-2. Intuitively preset $\zeta$ according to the decision-maker's (researcher's) desired level of risk aversion (Risk Preference)}

\begin{mdframed}[linewidth=0.6pt]
\footnotesize
\note \textbf{Kannan et al. (Residuals-based Distributionally Robust Optimization)}

\hrule\smallskip
\textbf{Algorithm 2} Specifying a covariate-independent radius $\zeta_n$ using the ER-DRO problem
\hrule\smallskip
\begin{enumerate}[label=\arabic*:,leftmargin=19pt,itemsep=2pt,parsep=0pt]
\item \textbf{Input:} data $\mathcal D_n$, set of candidate radii $\Delta$, number of folds $K$, and number of covariate realizations sampled during each fold $T\le\lfloor n/K\rfloor$.
\item Partition $[n]$ into subsets $S_1,\ldots,S_K$ of (roughly) equal size at random. Let $\mathcal D_{-k}:=\mathcal D_n\setminus\{(y^i,x^i)\}_{i\in S_k}$.
\item \textbf{for} $k=1,\ldots,K$ \textbf{do}
\item \hspace{6pt}Pick without replacement a random subset $\bar X$ of $\{x^i\}_{i\in S_k}$ of size $T$.
\item \hspace{6pt}\textbf{for} $\bar x\in\bar X$ \textbf{do}
\item \hspace{12pt}\textbf{for} $\zeta\in\Delta$ \textbf{do}
\item \hspace{18pt}Fit a regression model $\hat f_{-k}$ using the data $\mathcal D_{-k}$ and compute its in-sample residuals $\{\hat\varepsilon_{-k}^i\}_{i\notin S_k}:=\{y^i-\hat f_{-k}(x^i)\}_{i\notin S_k}$.
\item \hspace{18pt}Solve the ER-DRO problem below at covariate $\bar x$ to get solution $\hat z_{-k}^{\rm DRO}(\bar x,\zeta)$
\[
\min_{z\in Z}\sup_{Q\in\hat{\mathcal P}_{-k}(\bar x)}\mathbb E_{Y\sim Q}[c(z,Y)],
\]
where the ambiguity set $\hat{\mathcal P}_{-k}(\bar x)$ with radius $\zeta$ is centered at the estimated empirical distribution
\[
\hat P_{-k}^{\rm ER}(\bar x):=\frac1{n-|S_k|}\sum_{i\notin S_k}\delta_{\hat f_{-k}(\bar x)+\hat\varepsilon_{-k}^i}.
\]
\item \hspace{12pt}\textbf{end for}
\item \hspace{6pt}\textbf{end for}
\item \textbf{end for}
\item \textbf{Output:} Radius
\[
\zeta_n\in\arg\min_{\zeta\in\Delta}\frac1T\sum_{\bar x\in\bar X}\frac1K\sum_{k\in[K]}\frac1{|S_k|}\sum_{i\in S_k}c\bigl(\hat z_{-k}^{\rm DRO}(\bar x,\zeta),y^i\bigr)
\]
for the ambiguity set $\hat{\mathcal P}_n(x)$ for the ER-DRO problem (7).
\end{enumerate}
\hrule\smallskip
\textbf{Algorithm 3} Specifying a covariate-dependent radius $\zeta_n(x)$ using the ER-DRO problem
\hrule\smallskip
\begin{enumerate}[label=\arabic*:,leftmargin=19pt,itemsep=2pt,parsep=0pt]
\item \textbf{Input:} data $\mathcal D_n$, set of candidate radii $\Delta$, number of folds $K$, and new covariate realization $x\in X$.
\item Partition $[n]$ into subsets $S_1,\ldots,S_K$ of (roughly) equal size at random. Let $\mathcal D_{-k}:=\mathcal D_n\setminus\{(y^i,x^i)\}_{i\in S_k}$.
\item \textbf{for} $k=1,\ldots,K$ \textbf{do}
\item \hspace{6pt}\textbf{for} $\zeta\in\Delta$ \textbf{do}
\item \hspace{12pt}Fit a regression model $\hat f_{-k}$ using the data $\mathcal D_{-k}$ and compute its in-sample residuals $\{\hat\varepsilon_{-k}^i\}_{i\notin S_k}:=\{y^i-\hat f_{-k}(x^i)\}_{i\notin S_k}$.
\item \hspace{12pt}Solve the ER-DRO problem below at covariate $x$ to obtain solution $\hat z_{-k}^{\rm DRO}(x,\zeta)$
\[
\min_{z\in Z}\sup_{Q\in\hat{\mathcal P}_{-k}(x)}\mathbb E_{Y\sim Q}[c(z,Y)],
\]
where the ambiguity set $\hat{\mathcal P}_{-k}(x)$ with radius $\zeta$ is centered at the estimated empirical distribution
\[
\hat P_{-k}^{\rm ER}(x):=\frac1{n-|S_k|}\sum_{i\notin S_k}\delta_{\hat f_{-k}(x)+\hat\varepsilon_{-k}^i}.
\]
\item \hspace{12pt}Fit a regression model $\hat f_k$ using the data $\{(y^i,x^i)\}_{i\in S_k}$ and compute its in-sample residuals $\{\hat\varepsilon_k^i\}_{i\in S_k}:=\{y^i-\hat f_k(x^i)\}_{i\in S_k}$.
\item \hspace{6pt}\textbf{end for}
\item \textbf{end for}
\item \textbf{Output:} Radius
\[
\zeta_n(x)\in\arg\min_{\zeta\in\Delta}\frac1K\sum_{k\in[K]}\frac1{|S_k|}\sum_{i\in S_k}c\bigl(\hat z_{-k}^{\rm DRO}(x,\zeta),\hat f_k(x)+\hat\varepsilon_k^i\bigr)
\]
for the ambiguity set $\hat{\mathcal P}_n(x)$ for the ER-DRO problem (7).
\end{enumerate}
\hrule
\end{mdframed}
\end{document}
